\section{Training Procedure}\label{sec:method:training}

Parameters $\bm{\theta}$ are estimated by minimising $\Loss$
(\cref{eq:method:loss_total}) with Adam \parencite{kingma2014adam}.
Learning rate, weight decay, epoch count, dropout, and related
settings are listed in \cref{tab:exp:hparams_main}.

\subsection*{Scenario-aware batch sampling}\label{sec:method:training:sampler}

Mini-batches are not formed by uniform shuffling of the pooled windows.
Emitter identifiers are recording-local
(\cref{sec:method:formulation}), so two windows drawn from different
scenarios never share an emitter positive even when both contain
pulses from physically similar sources.
Uniform mixing across the corpus would therefore leave the emitter
branch with positives only inside each individual window, which is
often too few for a stable contrastive signal.
Each mini-batch is instead drawn from a single scenario, with several
windows taken from that scenario so that the same emitter can appear
across windows and supply additional positives.
Mode labels remain globally comparable, so the mode branch still sees
batch-wide positives and negatives under
\cref{sec:method:loss}.
The number of windows per batch, together with any early-stopping
policy, is recorded in \cref{tab:exp:hparams_main}.
